\section{Background and Motivation}
\subsection{The interval nobody owns}
Consider the reference site of this paper, scaledaiops.org: a static, volunteer-maintained reference for an AI-operations framework, without accounts or analytics. Until August 2026 its only feedback path was a link to the project's code-hosting organisation. A visitor without an account there could not report anything; one with an account had to leave the site, choose a repository and open an issue; nobody was notified; and no one could say how long such an issue typically waited. This is not unusual. Service desks and DevOps teams solve the same problem with people and process---ITIL-style acknowledgement, routing and resolution targets~\cite{axelos2019itil}, or flow metrics such as lead time and time to restore~\cite{forsgren2018accelerate}---but those presume staff whose job it is to work the queue.

\subsection{What ``small'' means here}
We use \emph{small project} operationally: at most three maintainers, no dedicated support or triage role, and fewer than one hundred feedback items per month. At that scale a vendor feedback product is disproportionate (data leaves the project, cost per seat, another login), a self-hosted tracker is another system to operate, and the raw Git host is unfriendly to outsiders. \ffrs{} targets exactly this gap.

\subsection{Why an agent changes the economics}
The scarce resource in a small project is maintainer attention. If the first substantive response to an item---a candidate fix with tests, or a concrete proposal---can be produced without it, attention is spent only on judgement: is the fix right, is the proposal acceptable. An independent commercial site that implemented a comparable capture widget (with a conventional, manual response process) reports the same pattern: capture works, items still stall. That observation motivated making the agentic stage a defining part of the model rather than an optional add-on.
The author's own path---systems administration in the 2000s, web development in the 2010s, then DevOps and release management for large delivery organisations---crossed the same failure repeatedly: feedback was captured diligently and acted on late or never, and nobody could say how late. What was missing was never a form; it was an owner for the interval after the form, and a clock on it. \ffrs{} is that owner and that clock, with the owner now an agent.
